\chapter{PART II · 思维重构}


\newpage

思维是行动的起点。

在商业世界里，思维通常不是被忽视的——恰恰相反，几乎每个组织都在强调"思维升级"、"认知迭代"、"心智模式转变"。但很多时候，这种强调只停留在口号层面。真正的思维转变，需要的是对旧思维的深刻理解，对新思维的清晰界定，以及在日常决策中持续实践的意志力。

这一部分要做的，不是提供一个新的思维框架，而是帮助读者识别那些在AI时代正在失效的思维模式，并为它们找到替代品。

四个章节，对应四种需要升级的思维：
\begin{itemize}
\item 从预测到塑造——认知范式的根本转移
\item 效率思维的陷阱——为什么做对的事比把事做对更重要
\item 最佳实践的半衰期——从复制已知到创造未知
\item 韧性优先——在不确定性中保持运转的能力

\end{itemize}
这四种思维陷阱，不是彼此独立的。它们相互关联，相互强化。效率思维让人追求最佳实践，最佳实践的失效让人更加依赖预测，而过度预测的组织往往缺乏韧性。打破这个循环，需要从任何一个环节入手——但最根本的，还是从认知范式的转移开始。
